\section{Related Work}
\label{sec:related-work}

This section positions the thesis against four strands of literature that map onto its
research questions: runway detection and the LARD benchmark, learning from synthetic
data and the Sim2Real gap, weather-robust augmentation, and domain adaptation via
self-training and generative style transfer. For each strand we make explicit how the
cited work relates to, and differs from, this thesis.

\subsection{Vision-Based Runway Detection and the LARD Benchmark}
\label{subsec:rw-runway}
Runway detection for vision-based landing has been approached both as bounding-box
detection and as geometry-aware localization. Runway-specific detectors such as
YOLO-RWY~\cite{yolorwy2024} train a dedicated YOLOv8 detection architecture on LARD-V1,
report performance on the real nominal and edge-case test splits, and introduce a
\emph{data-augmentation module} spanning diverse weather conditions (rain, snow, fog,
shadow, sun flare). Notably, the authors announce as future work the use of GANs for
data augmentation and feature-enhancement modules to compensate for the limited night
and adverse-weather data in LARD, as well as the extraction of runway
\emph{keypoints} for pose estimation --- directions this thesis takes up directly.
Beyond a single frame, landing-guidance systems combine detection with tracking and
orientation estimation: Ferreira et al.~\cite{ferreira2025landing} track the runway and
estimate aircraft orientation across an approach, and runway/airside safety systems use
detection-plus-tracking video analytics for surveillance~\cite{runwaysurveillance2026}.

The LARD benchmark~\cite{lard2023} provides both real and synthetic approach imagery
annotated with runway corners. Its second iteration, LARD-V2~\cite{syntheticrunway2025},
addresses two limitations of LARD-V1 --- an operational design domain (ODD) restricted
to single-runway airports with an unrealistic approach cone, and few generation sources
--- by narrowing the approach cone, extending to multi-runway airports, proposing an
extended-ODD metric (e-mAP) that accounts for both in-ODD and extended-ODD runways, and
expanding source diversity to improve cross-domain generalization. The authors benchmark
YOLOv11 detectors with mAP and e-mAP, training eleven models that differ only in
training-data composition (each individual source, the full dataset, and five
leave-one-out sets). Their findings are directly relevant here: the \emph{full}
configuration achieves the best average performance (every source contributes); the
satellite-only sources ArcGIS and Bing Maps combine well with each other but less well
with full 3D simulators; and \texttt{ges} attains the highest mAP among individual
sources, likely because it blends satellite imagery with 3D rendering.

\emph{Relation to this thesis.} Like YOLO-RWY, this thesis uses a lightweight YOLO
detector, but formulates the task as four-corner \emph{pose} estimation and focuses on
\emph{closing the Sim2Real gap} for a fixed detector rather than on architecture design.
Whereas the LARD-V2 authors studied cross-\emph{dataset} generalization with a single
detector on synthetic test sets, this thesis repeats the analysis against the
\emph{real} LARD-V1 set, stratifies it into weather scenarios as weather-specific
benchmarks (extended with private landing sequences to compensate for the sparse
adverse-weather representation), and evaluates weather augmentation, self-training and
style transfer as remedies. The strong YOLO benchmark and the object-detection findings,
which the LARD-V2 authors relate to the datasets' rendering foundations, transfer
naturally to the pose-estimation setting adopted here. Real-time detector alternatives
such as RT-DETR~\cite{rtdetr2024} are noted as possible baselines but are outside the
scope, which fixes the detector to isolate the adaptation strategy.

\subsection{Learning from Synthetic Data and the Sim2Real Gap}
\label{subsec:rw-sim2real}
Synthetic imagery from simulators and geospatial renderers enables large, exactly
labeled training sets, and domain randomization~\cite{tobin2017domainrand} broadens the
training distribution to improve transfer. For runways specifically, LARD-V2 generates
imagery from five pipelines (flight simulator, ArcGIS, Bing Maps, Google Earth Studio,
X-Plane) and shows that mixing sources improves cross-domain
generalization~\cite{syntheticrunway2025}. These works establish that synthetic data is
viable but that a residual Sim2Real gap remains, especially away from nominal
conditions.

\emph{Relation to this thesis.} This thesis first \emph{quantifies} that residual gap
in a controlled, shared-location cross-source comparison against the real domain
(\Cref{sec:gap}). Its methodological contribution is to show that the gap is only
measurable after controlling for a runway \emph{orientation} bias (vertical angle) in
addition to runway size --- without which the real set spuriously appears easier than
the synthetic sets.

\subsection{Weather-Robust Augmentation}
\label{subsec:rw-weather}
Detector accuracy degrades markedly under rain, fog, glare and low light, and
library-based augmentation such as \texttt{albumentations}~\cite{albumentations2020}
provides fast, composable weather-like transformations; robustness studies quantify how
much such degradations move detection metrics~\cite{weatherrobust2025}. Closest to this
thesis, Rüter et al.~\cite{ruter2025augmentations} systematically investigate how
individual image augmentations affect the sim-to-real generalization of perception
models, reporting per-augmentation performance changes relative to a baseline and
highlighting the combinatorial explosion of augmentation configurations.

\emph{Relation to this thesis.} \Cref{sec:weather-aug} applies weather augmentations to
the runway-pose task and, going beyond an on/off comparison, sweeps the augmentation
\emph{ratio} per weather scenario, controls the dominant sampling variance by fixing the
training set, and tests whether appearance augmentation can substitute for the diversity
of a multi-source training set.

\subsection{Self-Training, Pseudo-Labeling and Temporal Consistency}
\label{subsec:rw-selftraining}
Unsupervised domain adaptation aligns a model to an unlabeled target domain; classic
approaches range from feature alignment by backpropagation~\cite{ganin2015dann} to
self-training, where confident model predictions on target data serve as
pseudo-labels~\cite{lee2013pseudolabel}, typically in a student--teacher setup combining
pseudo-label and consistency objectives. In video, object trackers can propagate labels
across frames to enforce temporal consistency and recover objects the per-frame detector
misses~\cite{ferreira2025landing, runwaysurveillance2026, ostrack2022}.

\emph{Relation to this thesis.} \Cref{sec:pseudo-labels} instantiates confidence-filtered
self-training on unlabeled real landing sequences, compares thresholds, analyzes the
pseudo-label composition across iterations, and extends the loop with a tracker
(OSTrack) that propagates corner pseudo-labels forward and backward along a sequence to
improve their completeness for distant runways.

\subsection{Generative Style Transfer for Domain Adaptation}
\label{subsec:rw-style}
Image-to-image translation adapts source images to look like the target domain.
CycleGAN~\cite{zhu2017cyclegan} popularized unpaired translation for Sim2Real, but its
lack of explicit geometric constraints can hallucinate or displace structure and thus
break keypoint labels. Conditional diffusion models with a control branch, such as the
Cosmos platform~\cite{nvidia2025cosmos}, can be guided by the source structure (edge or
segmentation maps, reference images, prompts) to repaint appearance while preserving
scene geometry.

\emph{Relation to this thesis.} \Cref{sec:style-transfer} uses control-guided,
structure-preserving diffusion to import real and adverse-weather appearance onto the
combined synthetic source \emph{while keeping the four-corner labels valid}, and
optionally contrasts it with a CycleGAN baseline to quantify the label-displacement risk.
